% Generado por scripts/validar_modelo_financiero.py — no editar a mano.
\begin{table}[t]
  \centering
  \footnotesize\setlength{\tabcolsep}{3.5pt}
  \caption{Flujo de caja del escenario neutral, en USD. El flujo neto resta de los ingresos los costos, el impuesto a las ganancias y la inversión en capital de trabajo; en el año 5 esa inversión es negativa porque se recupera. Fuente: elaboración propia sobre la planilla del modelo financiero.}
  \label{tab:financiero-neutral-flujo}
  \begin{tabular}{c r r r r r r r}
    \hline
    \textbf{Año} & \textbf{Ingresos} & \textbf{Costos} & \textbf{Impuesto} & \textbf{Cap. de trabajo} & \textbf{Flujo neto} & \textbf{Descontado} & \textbf{Acumulado} \\
    \hline
    0 & 0,00 & 1.741,00 & 0,00 & 0,00 & (1.741,00) & (1.741,00) & (1.741,00) \\
    1 & 0,00 & 41.353,99 & 0,00 & 0,00 & (41.353,99) & (33.083,19) & (43.094,99) \\
    2 & 57.600,00 & 79.033,16 & 0,00 & 9.600,00 & (31.033,16) & (19.861,22) & (74.128,15) \\
    3 & 209.800,00 & 154.100,04 & 0,00 & 25.366,67 & 30.333,29 & 15.530,65 & (43.794,86) \\
    4 & 259.800,00 & 241.314,14 & 3.380,18 & 8.333,33 & 6.772,35 & 2.773,95 & (37.022,51) \\
    5 & 369.400,00 & 322.935,11 & 16.262,71 & (43.300,00) & 73.502,18 & 24.085,19 & 36.479,67 \\
    \hline
  \end{tabular}
\end{table}
